\section{从 Demo 到生产：人类仍然负责 Objective Function}

\subsection{人类的工作没有消失，只是后移到了目标定义层}
Karpathy 多次强调应当把人从 loop 中移除，但这并不等于人类不再重要。恰恰相反，人类的职责从 ``执行每一步'' 后移到了 ``定义 objective、设边界、提供 verifier''。Auto research 之所以能工作，不是因为 agent 会神奇地产生价值，而是因为目标函数足够清晰。

\begin{importantbox}{Agent 最怕的不是能力不足，而是目标函数模糊}
\begin{itemize}
    \item 写 CUDA kernel、更快通过 benchmark、降低训练 loss，这类任务可以用客观指标判定；
    \item 产品品味、用户意图、组织优先级，则往往没有单一标准答案；
    \item 一旦 objective function 不清晰，系统就会把大量 token 浪费在局部最优和错误方向上。
\end{itemize}
\end{importantbox}

这也是他为什么会把 ``I shouldn't be looking at the results'' 和 ``if you can't evaluate, you can't auto research it'' 放在一起说。前一句讨论的是去掉人工瓶颈，后一句讨论的是自动化的适用边界。两者组合起来，构成了一个非常工程化的判断：\textbf{凡是可验证的，就尽量放进闭环；凡是不可验证的，就不要假装它已经自动化。}

\subsection{从 agent-first API 到公司基础设施}
如果未来 ``apps shouldn't exist''，那软件公司最重要的资产就不再只是前端界面，而是其背后的 API、数据权限模型、可观测性和审计日志。Agent 要成为主要使用者，系统就必须天然暴露给 agent，而不是靠 UI 自动化去逆向操作。

\begin{knowledgebox}{Agent-first 产品的基础设施清单}
\begin{enumerate}
    \item \textbf{可组合 API}：稳定、可授权、可审计；
    \item \textbf{状态与记忆层}：让 agent 能跨会话延续任务；
    \item \textbf{可观测性}：知道 agent 做了什么、为什么失败；
    \item \textbf{回滚与权限控制}：避免长链路自动化把错误放大。
\end{enumerate}
\end{knowledgebox}

\begin{table}[H]
\centering
\begin{tabular}{p{3.0cm}p{4.1cm}p{4.8cm}}
\toprule
\textbf{问题} & \textbf{旧的软件假设} & \textbf{agent-first 假设} \\
\midrule
主要用户是谁 & 人类点击 UI & agent 调用 API，人类做监督与验收 \\
产品入口是什么 & 页面、按钮、表单 & message channel、API、workflow trigger \\
如何提升效率 & 简化交互流程 & 暴露操作能力、减少上下文切换 \\
如何积累壁垒 & 更精致的前端与流程 & 数据、工具接入、权限模型与 evaluation loop \\
\bottomrule
\end{tabular}
\caption{Karpathy 所说的 agent-first，并不是多一个聊天框，而是软件分发逻辑的重写}
\end{table}

\subsection{本章小结}
人类不会从系统中消失，但会更多站在目标函数、权限边界和验收标准的位置上。真正的生产化重点，是把 agent 能力接进 API、状态管理和可观测性，而不是停留在 demo 式自动化。


